\documentclass[11pt,a4paper]{article}
\usepackage{shared/luxcover}
\usepackage[utf8]{inputenc}
\usepackage[T1]{fontenc}
\usepackage{amsmath,amssymb,amsfonts,amsthm}
\usepackage{graphicx}
\usepackage{hyperref}
\usepackage{algorithm}
\usepackage{algpseudocode}
\usepackage{booktabs}
\usepackage{listings}
\input{shared/lstlang}
\usepackage{xcolor}
\usepackage{enumitem}
\usepackage{float}

\newtheorem{theorem}{Theorem}
\newtheorem{lemma}[theorem]{Lemma}
\newtheorem{definition}{Definition}

\lstset{
    basicstyle=\small\ttfamily,
    keywordstyle=\color{blue}\bfseries,
    commentstyle=\color{gray},
    stringstyle=\color{red},
    breaklines=true,
    frame=single,
    numbers=left,
    numberstyle=\tiny\color{gray}
}

\hypersetup{
    colorlinks=true,
    linkcolor=blue,
    citecolor=blue,
    urlcolor=blue
}

\title{Pasta Curves Precompile:\\Pallas and Vesta for Halo2 Proof Verification}
\author{
Zach Kelling\thanks{Corresponding author: zach@lux.network}\\
\textit{Lux Industries, Inc.}\\
\texttt{research@lux.network}
}
\date{April 2026}

\begin{document}

\luxcoverpage

\begin{abstract}
We specify the Pasta curves precompile deployed on the Lux EVM, implementing
point addition, scalar multiplication, and multi-scalar multiplication over both
the Pallas and Vesta elliptic curves. The Pasta cycle (Pallas And Vesta)
is a pair of prime-order curves where each curve's scalar field equals the other's
base field, forming a 2-cycle essential for recursive proof composition in the
Halo2 proving system. The precompile enables on-chain verification of
Halo2 proofs (used in Zcash, Scroll, and Taiko), IPA (inner product argument)
commitment verification, and recursive SNARK aggregation. Gas costs are 150 gas
for addition and 7\,000 gas for scalar multiplication on either curve.
\end{abstract}

\section{Introduction}

The Halo2 proof system~\cite{halo2} uses inner product arguments (IPA) over
Pedersen commitments rather than pairings, eliminating the trusted setup required
by Groth16. IPA verification requires multi-scalar multiplication over a
commitment key, which is prohibitively expensive in EVM bytecode.

The Pasta curves~\cite{pasta} are specifically designed for Halo2:
\begin{enumerate}[label=(\roman*)]
  \item \textbf{Pallas}: base field $\mathbb{F}_p$, scalar field $\mathbb{F}_q$
  \item \textbf{Vesta}: base field $\mathbb{F}_q$, scalar field $\mathbb{F}_p$
\end{enumerate}
where $p$ and $q$ are 255-bit primes. The 2-cycle property ($p_{\text{Pallas}} = q_{\text{Vesta}}$
and $q_{\text{Pallas}} = p_{\text{Vesta}}$) enables recursive proof composition:
a Pallas circuit can verify a Vesta proof natively, and vice versa.

\section{Curve Parameters}

\begin{definition}[Pallas]
Short Weierstrass form $y^2 = x^3 + 5$ over $\mathbb{F}_p$ with:
\begin{align}
  p &= \texttt{0x40000000000000000000000000000000224698fc094cf91b992d30ed00000001} \\
  q &= \texttt{0x40000000000000000000000000000000224698fc0994a8dd8c46eb2100000001}
\end{align}
The curve has prime order $q$ (cofactor 1).
\end{definition}

\begin{definition}[Vesta]
Short Weierstrass form $y^2 = x^3 + 5$ over $\mathbb{F}_q$ with:
\begin{align}
  \text{base field} &= q \quad \text{(Pallas scalar field)} \\
  \text{scalar field} &= p \quad \text{(Pallas base field)}
\end{align}
Also prime order (cofactor 1).
\end{definition}

\begin{theorem}[2-Cycle]
Pallas and Vesta form a 2-cycle: arithmetic in Pallas's scalar field is
arithmetic in Vesta's base field. A circuit over Pallas can perform native
field arithmetic on Vesta scalars without range checks or non-native arithmetic.
\end{theorem}

\section{Precompile Specification}

\subsection{Address Assignment}

Address: \texttt{0x0200\ldots22}. Selector byte encodes curve and operation:

\begin{table}[H]
\centering
\begin{tabular}{lllr}
\toprule
\textbf{Selector} & \textbf{Curve} & \textbf{Operation} & \textbf{Gas} \\
\midrule
\texttt{0x01} & Pallas & Point Addition & 150 \\
\texttt{0x02} & Pallas & Scalar Multiplication & 7\,000 \\
\texttt{0x03} & Pallas & MSM & 7\,000$+$ \\
\texttt{0x11} & Vesta & Point Addition & 150 \\
\texttt{0x12} & Vesta & Scalar Multiplication & 7\,000 \\
\texttt{0x13} & Vesta & MSM & 7\,000$+$ \\
\bottomrule
\end{tabular}
\end{table}

\subsection{Point Encoding}

Points are encoded in uncompressed affine form: 32-byte $x$ followed by 32-byte $y$,
big-endian. The identity point is 64 zero bytes.

\subsection{Multi-Scalar Multiplication}

MSM is critical for IPA verification. A Halo2 proof verification requires one MSM
of size $2n$ where $n$ is the number of IPA rounds (typically 10--15, giving
MSM of 1024--32768 points).

The implementation uses the Pippenger bucket method with optimal window size
$c = \lceil \log_2 n \rceil$. The gas formula matches BLS12-381 MSM with the
EIP-2537 discount schedule, scaled to the Pasta field size.

\section{Security Analysis}

\subsection{Prime Order}

Both Pallas and Vesta have cofactor 1 (prime order). There is no small subgroup
attack. No cofactor multiplication is needed. This is the simplest possible
security model for an elliptic curve precompile.

\subsection{Field Validation}

Input coordinates are checked to be in $[0, p)$ (Pallas) or $[0, q)$ (Vesta).
Points are verified on-curve: $y^2 \equiv x^3 + 5$. Invalid inputs return error.

\subsection{Curve Confusion}

The selector byte distinguishes Pallas from Vesta operations. The precompile
validates that input coordinates are in the correct field for the selected curve.
A Pallas point submitted to a Vesta operation (or vice versa) is rejected if the
coordinates exceed the target field modulus.

\section{Gas Cost Justification}

Pallas and Vesta use 255-bit non-Mersenne primes, slightly slower than
Curve25519 ($2^{255}-19$) but comparable to BN254 base field. The cofactor-1
property eliminates cofactor clearing costs. Gas costs are between Curve25519
(5\,000 for ScalarMul) and BLS12-381 G1 (12\,000 for ScalarMul).

\begin{table}[H]
\centering
\begin{tabular}{lrr}
\toprule
\textbf{Operation} & \textbf{Time ($\mu$s)} & \textbf{Gas} \\
\midrule
Pallas Add & 2.3 & 150 \\
Pallas ScalarMul & 108 & 7\,000 \\
Pallas MSM (256) & 6\,400 & 280\,000 \\
Vesta Add & 2.3 & 150 \\
Vesta ScalarMul & 108 & 7\,000 \\
\bottomrule
\end{tabular}
\end{table}

\section{Halo2 Verification}

A complete Halo2 proof verification on-chain requires:
\begin{enumerate}
  \item Commitment opening check: 1 MSM of size $2^k$ ($k$ = circuit size parameter).
  \item Permutation argument: $O(1)$ scalar multiplications.
  \item Lookup argument: $O(1)$ scalar multiplications.
\end{enumerate}

With the MSM precompile, a Halo2 proof over a circuit with $2^{20}$ rows
can be verified in $\sim$5M gas, compared to $\sim$500M gas with Solidity-only MSM.

\section{Conclusion}

The Pasta curves precompile enables on-chain Halo2 proof verification at
practical gas costs. The 2-cycle property is preserved through careful
selector-based dispatch, and the prime-order curves eliminate cofactor
complications. The MSM operation with Pippenger optimization makes IPA
verification 100x cheaper than pure Solidity.

\begin{thebibliography}{9}
\bibitem{halo2}
  S.~Bowe, J.~Grigg, and D.~Hopwood, ``Recursive Proof Composition without
  a Trusted Setup,'' 2019. Revised in Halo2 book, 2022.

\bibitem{pasta}
  D.~Hopwood, ``The Pasta Curves for Halo 2 and Beyond,'' Zcash specification, 2020.
\end{thebibliography}

\end{document}
